\scanpage{105}
\uline{Théorème} (Seshadri) Soit $f:X\longrightarrow S$ un morphisme propre, plat, \deleted{\unsure{$(S_2)$}}, tel que pour $\forall s\in S$, \deleted{\unsure{on ait} $k(s)\simeq H^0(X_s,\mathcal O_{X_s})$ et} $X_s$ soit $(S_2)$ \inserted{(\unsure{et séparable} $(R_0)$)}. On suppose $S$ loc. noeth normal irréductible de pt gén. $\eta$, \deleted{et que $X/S$ est localement \unsure{pourvu} \inserted{\unsure{par} \unsure{extension}} de sections étales}. On suppose que $\underline{\operatorname{Pic}}_{X/S}$ existe, et est \uline{essentiellement propre} sur $S$. [ce dernier fait étant vérifié si $f$ est simple, ou $X=S\times_kY$, $Y$ schéma normal$/k$, $k$ un corps].

Soit $D$ un \deleted{diviseur} \inserted{$1$-cycle} sur $X$, tel que

a) $\operatorname{supp}D=Z$ ne contient \unsure{aucune} \deleted{\unsure{$\eta$}} \deleted{d'aucune} \inserted{une} fibre de $X$, i.e. $\Leftarrow Z_s$ est \uline{maigre} dans $X_s$.

b) $D_\eta$ est un diviseur de Cartier.

\uline{Alors $D$ est un diviseur \deleted{de Cartier}}.

\margin{\uline{N.B.} Le raisonnement direct sur $\operatorname{Div}^+_{X/S}$ marche seulement si $D\geq0$ !}
\margin{\deleted{\unsure{On peut} \unsure{raisonner} \unsure{directement} \unsure{sur} $\underline{\operatorname{Div}}^+$.}}

On va \uline{d'abord} supposer que $X/S$ a une section $\sigma$. Notons que $D_\eta$ définit un faisceau inversible $\underline L_\eta$, donc une section \inserted{rationnelle} de $\underline{\operatorname{Pic}}_{X/S}$ sur $S$, d'où un Module
\margin{104}
